  \documentclass{report}

% Language setting
% Replace `english' with e.g. `spanish' to change the document language
\usepackage[utf8]{inputenc}
\usepackage[polish]{babel}
\usepackage[T1]{fontenc}

\usepackage{array}
\usepackage{booktabs}

% Set page size and margins
% Replace `letterpaper' with `a4paper' for UK/EU standard size
\usepackage[letterpaper,top=2cm,bottom=2cm,left=3cm,right=3cm,marginparwidth=1.75cm]{geometry}

% Useful packages
\usepackage{amsmath}
\usepackage{graphicx}
\usepackage[colorlinks=true, allcolors=black]{hyperref}

\usepackage{float}


\usepackage{listings}
\lstset{
  inputencoding=utf8,
  extendedchars=true,
  literate=
    {ą}{{\k{a}}}1
    {ć}{{\'c}}1
    {ę}{{\k{e}}}1
    {ł}{{\l{}}}1
    {ń}{{\'n}}1
    {ó}{{\'o}}1
    {ś}{{\'s}}1
    {ź}{{\'z}}1
    {ż}{{\.z}}1
    {Ą}{{\k{A}}}1
    {Ć}{{\'C}}1
    {Ę}{{\k{E}}}1
    {Ł}{{\L{}}}1
    {Ń}{{\'N}}1
    {Ó}{{\'O}}1
    {Ś}{{\'S}}1
    {Ź}{{\'Z}}1
    {Ż}{{\.Z}}1,
}

\usepackage{xcolor}

\definecolor{codegreen}{rgb}{0,0.6,0}
\definecolor{codegray}{rgb}{0.5,0.5,0.5}
\definecolor{codepurple}{rgb}{0.58,0,0.82}
\definecolor{backcolour}{rgb}{0.95,0.95,0.92}

\lstdefinestyle{mystyle}{
    backgroundcolor=\color{backcolour},   
    commentstyle=\color{codegreen},
    keywordstyle=\color{magenta},
    numberstyle=\tiny\color{codegray},
    stringstyle=\color{codepurple},
    basicstyle=\ttfamily\footnotesize,
    breakatwhitespace=false,         
    breaklines=true,                 
    captionpos=b,                    
    keepspaces=true,                 
    numbers=left,                    
    numbersep=5pt,                  
    showspaces=false,                
    showstringspaces=false,
    showtabs=false,                  
    tabsize=2
}

\lstset{style=mystyle}

\usepackage{tikz}
\usetikzlibrary{trees}
\usepackage{forest}
\usetikzlibrary{arrows.meta, shapes.geometric, positioning}

\usepackage{authblk}

\usepackage{enumitem}
\usepackage{float}
\usepackage{subcaption}

\title{\LARGE{Analiza skupień danych dotyczących cen mieszkań}

\vspace{0.3cm}

\Large{Metody Analizy Danych - zadanie projektowe 1}}

\author{
Joanna Dagil (231008),
Liliana Grześkiewicz (236138),
Piotr Szpatusko (223420),
Adam Włodarski (230983)
}

\date{rok akademicki 2025/2026}

\begin{document}

\maketitle

\tableofcontents

\chapter*{Streszczenie}
\addcontentsline{toc}{chapter}{Streszczenie}

Celem pracy była identyfikacja jednorodnych grup ofert sprzedaży mieszkań w Krakowie z wykorzystaniem metod analizy skupień. 
Materiał badawczy stanowiły dane pochodzące ze zbioru Kaggle obejmującego ogłoszenia mieszkaniowe z 15 największych miast Polski, przy czym analizę zawężono do rynku krakowskiego oraz wyłącznie do ofert sprzedaży. 
W badaniu uwzględniono cechy opisujące mieszkania i ich otoczenie, takie jak cena, powierzchnia, liczba pokoi, liczba punktów zainteresowania w okolicy, odległość od centrum oraz informacja o posiadaniu balkonu. 
Po wstępnym przygotowaniu danych, obejmującym ocenę kompletności danych, usunięcie obserwacji odstających oraz standaryzację zmiennych, zastosowano algorytm K-means oraz hierarchiczną metodę Warda. 
Uzyskane wyniki wskazują na istnienie trzech interpretowalnych segmentów rynku mieszkaniowego, różniących się poziomem cen, wielkością lokali oraz położeniem względem centrum miasta. 
Zbieżność wyników obu metod potwierdza wewnętrzne zróżnicowanie krakowskiego rynku mieszkań.

\vspace{0.5cm}
\noindent\textbf{Słowa kluczowe:} analiza skupień, klastry, ceny mieszkań, K-means, metoda Warda

\chapter{Wstęp}

Rynek mieszkaniowy należy do obszarów charakteryzujących się dużym zróżnicowaniem ofert pod względem ceny, wielkości lokalu, układu funkcjonalnego oraz położenia. 
W praktyce oznacza to, że mieszkania dostępne na tym samym rynku lokalnym mogą znacząco różnić się między sobą, tworząc naturalne segmenty odpowiadające odmiennym potrzebom nabywców oraz różnym poziomom dostępności cenowej. 
Z tego względu analiza rynku nieruchomości wymaga narzędzi pozwalających uchwycić nie tylko pojedyncze zależności między zmiennymi, lecz także wewnętrzną strukturę całego zbioru danych.

Jednym z podejść umożliwiających takie badanie jest analiza skupień, zaliczana do metod uczenia nienadzorowanego. 
Jej celem jest grupowanie obiektów podobnych do siebie pod względem wybranych cech, bez konieczności wcześniejszego przypisywania im etykiet. 
W kontekście rynku mieszkaniowego pozwala to na identyfikację segmentów ofert o zbliżonych parametrach, a następnie na ich interpretację ekonomiczną.

Podstawę analizy stanowi zbiór danych pozyskany z serwisu Kaggle, obejmujący oferty sprzedaży oraz najmu mieszkań z 15 największych miast w Polsce.
W toku prac badanie zawężono do ofert sprzedaży mieszkań w Krakowie. 
Wybór tego miasta wynikał z dużej liczby dostępnych obserwacji, co umożliwia bardziej wiarygodne zastosowanie metod grupowania oraz lepsze uchwycenie zróżnicowania badanego rynku. 
Ograniczenie analizy do ofert sprzedaży pozwoliło dodatkowo zachować jednorodność ekonomiczną danych i uniknąć mieszania informacji o cenach sprzedaży z cenami najmu.

\textbf{Głównym celem projektu} jest identyfikacja jednorodnych grup ofert sprzedaży mieszkań w Krakowie na podstawie cech opisujących lokale oraz ich otoczenie. 
W analizie wykorzystano zmienne odnoszące się do ceny mieszkania, powierzchni, liczby pokoi, dostępności punktów zainteresowania, odległości od centrum miasta oraz informacji o posiadaniu balkonu. 
Do realizacji celu zastosowano dwie metody analizy skupień: algorytm K-means oraz hierarchiczną metodę Warda.

Przeprowadzone badanie ma charakter eksploracyjny. 
Jego celem nie jest prognozowanie cen ani ocena pojedynczych nieruchomości, lecz rozpoznanie struktury badanego rynku i wskazanie, czy wśród ofert sprzedaży mieszkań w Krakowie można wyróżnić wyraźne, interpretowalne segmenty. 
Wyniki takiej analizy mogą pomóc w lepszym zrozumieniu wewnętrznego zróżnicowania rynku mieszkaniowego oraz stanowić punkt wyjścia do dalszych badań.

Realizacja projektu ma na celu:
\begin{itemize}
    \item lepsze zrozumienie struktury polskiego rynku mieszkań;
    \item identyfikację kluczowych segmentów rynku różniących się poziomem cen;
    \item stworzenie solidnej podstawy do dalszej analizy trendów cenowych oraz porównań regionalnych.
\end{itemize}

\chapter{Przedmiot badania}

Przedmiotem badania są oferty mieszkań pochodzące ze zbioru danych udostępnionego w serwisie Kaggle, obejmującego ogłoszenia sprzedaży i najmu z 15 największych miast Polski. Na potrzeby niniejszej pracy analizę zawężono do jednego miasta, aby uzyskać bardziej jednorodny i porównywalny zbiór obserwacji, umożliwiający przeprowadzenie pogłębionej analizy skupień.

W pierwszym etapie porównano ogólną liczbę dostępnych ofert w poszczególnych miastach. Na podstawie tej analizy do dalszego badania wybrano Kraków, który charakteryzuje się bardzo dużą liczbą ogłoszeń, ustępując pod tym względem jedynie Warszawie. Duża liczebność obserwacji zwiększa wiarygodność wyników i pozwala lepiej uchwycić zróżnicowanie badanego rynku.

\begin{figure}[H]
    \centering
    \includegraphics[width=0.99\linewidth]{miasta_najwiecej_wynajmow.png}
    \caption{Liczba ofert mieszkaniowych w poszczególnych miastach z podziałem na miesiące}
\end{figure}

Choć zbiór danych zawiera zarówno oferty sprzedaży, jak i najmu, dalszą analizę ograniczono wyłącznie do ofert sprzedaży mieszkań w Krakowie. Decyzja ta wynika z faktu, że ceny sprzedaży i ceny najmu opisują różne zjawiska ekonomiczne, przez co ich łączne analizowanie mogłoby prowadzić do niejednoznacznych i trudnych do interpretacji wyników. Ograniczenie badania do jednego typu ofert pozwala zachować spójność analizowanego zbioru.

Jednostką badania jest pojedyncza oferta mieszkania opisana za pomocą zestawu cech ilościowych i jakościowych. W analizie uwzględniono zmienne odnoszące się zarówno do właściwości samego lokalu, jak i do jego położenia oraz otoczenia. Są to: cena mieszkania, powierzchnia użytkowa, liczba pokoi, liczba punktów zainteresowania w promieniu 500 metrów, odległość od centrum miasta oraz informacja o posiadaniu balkonu.

W celu zapewnienia jakości dalszych analiz zbiór danych został poddany wstępnemu przygotowaniu, obejmującemu wybór ofert sprzedaży z Krakowa, ocenę kompletności danych oraz usunięcie obserwacji odstających. Tak przygotowany materiał badawczy stanowi podstawę do zastosowania metod analizy skupień, których celem jest identyfikacja grup mieszkań o podobnej charakterystyce.

\section{Cel i zakres badania}

Celem badania jest identyfikacja jednorodnych grup ofert mieszkań na rynku krakowskim przy wykorzystaniu metod analizy skupień. W szczególności chodzi o sprawdzenie, czy na podstawie wybranych cech nieruchomości możliwe jest wyodrębnienie segmentów rynku różniących się poziomem cen, wielkością mieszkań, liczbą pokoi, dostępnością infrastruktury w otoczeniu oraz położeniem względem centrum miasta.

Badanie ma charakter eksploracyjny i wykorzystuje metody uczenia nienadzorowanego, co oznacza, że analizowane dane nie zawierają z góry przypisanych kategorii mieszkań. Zamiast tego celem jest wykrycie naturalnej struktury zbioru danych i wskazanie, jakie typy ofert występują na badanym rynku. Pozwala to spojrzeć na rynek nieruchomości nie przez pryzmat pojedynczych ogłoszeń, lecz przez pryzmat powtarzalnych segmentów o podobnych właściwościach.

Zakres badania został zawężony do ofert dotyczących miasta Krakowa. Wybór tego miasta wynika z dużej liczby dostępnych obserwacji, co zwiększa reprezentatywność wyników i pozwala na bardziej wiarygodne zastosowanie metod grupowania. Dalszą analizę ograniczono wyłącznie do ofert sprzedaży mieszkań, ponieważ łączne analizowanie cen sprzedaży i cen najmu prowadziłoby do nieporównywalności obserwacji oraz mogłoby zaburzać wyniki analizy skupień.

W badaniu uwzględniono zmienne opisujące zarówno cechy fizyczne mieszkań, jak i ich otoczenie, tj. cenę, powierzchnię, liczbę pokoi, liczbę punktów zainteresowania w promieniu 500 metrów, odległość od centrum miasta oraz informację o posiadaniu balkonu. Przed właściwą analizą przeprowadzono wstępne przygotowanie danych obejmujące ograniczenie zbioru do ofert sprzedaży, ocenę kompletności danych oraz usunięcie obserwacji odstających.

W części analitycznej zastosowano dwie metody grupowania: algorytm K-means oraz hierarchiczną analizę skupień metodą Warda. Pozwoliło to nie tylko na wyznaczenie segmentów rynku, lecz także na porównanie wyników uzyskanych różnymi podejściami i ocenę ich spójności.

Realizacja celu badania będzie oceniana poprzez odpowiedź na trzy pytania: czy w danych można wskazać wyraźne skupienia ofert mieszkań, czy wyodrębnione klastry dają się sensownie zinterpretować ekonomicznie oraz czy wyniki uzyskane różnymi metodami analizy skupień są ze sobą zgodne. Tak sformułowany cel pozwala w końcowym podsumowaniu jednoznacznie ocenić, w jakim stopniu projekt doprowadził do poznania struktury badanego fragmentu rynku nieruchomości.





\section{Przegląd literatury}

%(min. 1 cytowanie powiązane tematycznie i krótki opis co było badane)

W literaturze dotyczącej rynku nieruchomości podkreśla się, że rynek mieszkaniowy nie ma charakteru jednorodnego, lecz składa się z wielu segmentów lub subrynków różniących się cenami, lokalizacją, standardem oraz cechami samych nieruchomości. Oznacza to, że mieszkania nie konkurują ze sobą w jednakowym stopniu, a podobieństwo ofert zależy zarówno od cech fizycznych lokalu, jak i od jego otoczenia. W konsekwencji badanie struktury rynku wymaga stosowania metod, które pozwalają identyfikować grupy obiektów podobnych do siebie wewnętrznie, a zarazem odmiennych od pozostałych.

Jedną z ważniejszych prac w tym obszarze jest artykuł Bourassy, Hamelinka, Hoesliego i MacGregora \cite{bourassa1999}. Autorzy zaproponowali statystyczny sposób wyodrębniania subrynków mieszkaniowych z wykorzystaniem analizy czynnikowej oraz analizy skupień. Następnie porównali jakość tak wyznaczonych podziałów z innymi sposobami klasyfikacji rynku. Wyniki ich badania pokazały, że empirycznie wyodrębnione subrynki mogą lepiej odzwierciedlać zróżnicowanie rynku niż podziały przyjęte arbitralnie. Praca ta stanowi istotne uzasadnienie dla stosowania metod grupowania w analizie danych mieszkaniowych.

W odniesieniu do warunków polskich szczególnie istotna jest praca Mateusza Tomala \cite{tomal2021}. Autor badał zróżnicowanie rynków mieszkaniowych w Polsce i wykazał, że można wyróżnić jednorodne klastry rynków różniących się wielkością zasobu mieszkaniowego, jego jakością oraz poziomem cen. Badanie to potwierdza, że także w Polsce rynek nieruchomości cechuje się wyraźną heterogenicznością, a zastosowanie metod klasteryzacji pozwala uchwycić jego strukturę w sposób uporządkowany i interpretowalny. Z perspektywy niniejszej pracy wniosek ten jest szczególnie ważny, ponieważ wspiera zasadność poszukiwania segmentów również na poziomie pojedynczego miasta.

Uzupełnieniem tego podejścia jest praca Tomala i Helbicha \cite{tomalhelbich2022}, poświęcona analizie subrynków mieszkaniowych w Krakowie. Autorzy wykazali, że nawet w obrębie jednego miasta możliwe jest wyróżnienie obszarów i grup ofert o odmiennej charakterystyce, a znaczenie poszczególnych cech mieszkań i ich lokalizacji może zmieniać się w czasie. Choć badanie dotyczyło rynku najmu, a nie sprzedaży, stanowi ono ważny punkt odniesienia dla niniejszego projektu, ponieważ potwierdza istnienie wewnętrznego zróżnicowania krakowskiego rynku mieszkaniowego oraz znaczenie cech lokalizacyjnych i otoczenia dla kształtowania się segmentów rynku.

Przywołane badania wskazują więc, że analiza skupień jest uzasadnionym narzędziem badania rynku mieszkaniowego. Literatura dostarcza zarówno podstaw teoretycznych dla pojęcia subrynku mieszkaniowego, jak i empirycznych przykładów pokazujących, że segmentacja rynku może być skutecznie analizowana metodami statystycznymi. Na tym tle niniejsza praca wpisuje się w istniejący nurt badań, koncentrując się na identyfikacji skupień ofert sprzedaży mieszkań w Krakowie na podstawie cech opisujących lokal, jego cenę oraz położenie.




\section{Zmienne wybrane do analizy}

W analizie przyjęto podział zmiennych diagnostycznych na symulanty i destymulanty z punktu widzenia atrakcyjności mieszkania. Symulantami nazywa się cechy, których wyższe wartości są oceniane pozytywnie, natomiast destymulantami są cechy, których wzrost wpływa niekorzystnie na ocenę obiektu. Dobór zmiennych miał na celu uwzględnienie zarówno cech fizycznych mieszkania, jak i elementów związanych z jego lokalizacją oraz dostępnością infrastruktury w otoczeniu.

\medskip

Zmienne wybrane do analizy to:

\begin{enumerate}[leftmargin=*, label=\arabic*.]
    \item[] \textbf{Symulanty}

    \item \texttt{squareMeters} -- powierzchnia mieszkania wyrażona w metrach kwadratowych

    Zmienna ta opisuje wielkość lokalu mieszkalnego i należy do podstawowych cech charakteryzujących nieruchomość. Powierzchnia mieszkania wpływa bezpośrednio na komfort użytkowania, możliwości aranżacyjne oraz funkcjonalność przestrzeni. Większy metraż daje zazwyczaj większą swobodę zagospodarowania wnętrza i lepiej odpowiada na potrzeby gospodarstw domowych o większej liczbie członków. Z tego względu zmienna \texttt{squareMeters} została zakwalifikowana jako symulanta.

    \item \texttt{rooms} -- liczba pokoi w mieszkaniu

    Zmienna określa liczbę wydzielonych pomieszczeń mieszkalnych w lokalu, a więc w pewnym stopniu odzwierciedla jego układ funkcjonalny. Liczba pokoi ma istotne znaczenie z punktu widzenia użyteczności mieszkania, ponieważ wpływa na możliwość oddzielenia poszczególnych stref życia codziennego, takich jak wypoczynek, sen czy praca. Mieszkania z większą liczbą pokoi są często bardziej atrakcyjne dla rodzin oraz osób pracujących zdalnie. Dlatego zmienna \texttt{rooms} została uznana za symulantę.

    \item \texttt{poiCount} -- liczba punktów zainteresowania w promieniu 500 metrów od mieszkania

    Zmienna ta opisuje intensywność i dostępność infrastruktury w bezpośrednim otoczeniu nieruchomości. Obejmuje ona liczbę punktów zainteresowania znajdujących się w promieniu 500 metrów od mieszkania, takich jak sklepy, punkty usługowe, restauracje, placówki edukacyjne, przystanki komunikacji miejskiej czy inne obiekty użyteczności codziennej. Wyższa wartość tej zmiennej wskazuje na lepsze wyposażenie okolicy w usługi i udogodnienia, co zwiększa wygodę życia mieszkańców oraz atrakcyjność lokalizacji. W związku z tym \texttt{poiCount} została zakwalifikowana jako symulanta.

    \item \texttt{hasBalcony} -- informacja o obecności balkonu w mieszkaniu

    Zmienna ta ma charakter jakościowy i informuje, czy mieszkanie posiada balkon. Balkon stanowi dodatkową przestrzeń użytkową, która może pełnić funkcję rekreacyjną, wypoczynkową lub praktyczną. Obecność balkonu jest często postrzegana jako cecha podnosząca standard mieszkania, szczególnie w zabudowie wielorodzinnej. Na potrzeby analizy zmienna ta może zostać zakodowana binarnie, na przykład jako 1 dla odpowiedzi \texttt{yes} oraz 0 dla odpowiedzi \texttt{no}. Po takim przekształceniu wyższa wartość oznacza cechę bardziej pożądaną, dlatego \texttt{hasBalcony} potraktowano jako symulantę.

    \item[] \textbf{Destymulanty}

    \item \texttt{price} -- cena mieszkania

    Zmienna określa całkowitą cenę oferowanej nieruchomości i stanowi jedną z najważniejszych cech z punktu widzenia nabywcy. Cena wpływa bezpośrednio na dostępność ekonomiczną mieszkania oraz możliwość jego zakupu przez potencjalnych klientów. Przy porównywaniu mieszkań wyższa cena oznacza zwykle mniejszą dostępność finansową, nawet jeśli może częściowo wynikać z lepszej lokalizacji lub wyższego standardu. Z tego względu, w kontekście oceny atrakcyjności mieszkania z perspektywy ekonomicznej, zmienna \texttt{price} została uznana za destymulantę.

    \item \texttt{centreDistance} -- odległość mieszkania od centrum miasta w kilometrach

    Zmienna ta opisuje położenie nieruchomości względem centralnej części miasta. Lokalizacja jest jednym z kluczowych czynników wpływających na wartość i atrakcyjność mieszkania, ponieważ determinuje dostęp do miejsc pracy, usług, transportu publicznego, oferty handlowej i kulturalnej. W warunkach dużego miasta mniejsza odległość od centrum jest zazwyczaj postrzegana jako cecha korzystna, gdyż wiąże się z lepszą dostępnością najważniejszych funkcji miejskich. Wobec tego wzrost wartości zmiennej \texttt{centreDistance} obniża atrakcyjność nieruchomości, dlatego została ona zakwalifikowana jako destymulanta.

\end{enumerate}

\section{Wstępna analiza danych}

\subsection{Statystyki opisowe i podstawowa wizualizacja}

% (przynajmniej: średnia, mediana, minimum, maksimum, odchylenie standardowe, skośność)

\begin{table}[H]

\caption{Statystyki opisowe zmiennych wykorzystanych w analizie}
\centering
\begin{tabular}[t]{lrrrrrr}
\toprule
  & Srednia & Mediana & Minimum & Maksimum & OdchylenieStandardowe & Skosnosc\\
\midrule
price & 626 200.24 & 660 000.00 & 1 700.00 & 3 000 000.00 & 526 704.50 & 0.75\\
squareMeters & 55.35 & 50.49 & 25.00 & 150.00 & 21.38 & 1.48\\
rooms & 2.53 & 2.00 & 1.00 & 6.00 & 0.92 & 0.99\\
poiCount & 24.35 & 14.00 & 0.00 & 212.00 & 33.62 & 2.99\\
centreDistance & 3.82 & 3.71 & 0.14 & 9.99 & 2.17 & 0.38\\
hasBalcony & 0.63 & 1.00 & 0.00 & 1.00 & 0.48 & -0.55\\
\bottomrule
\end{tabular}
\end{table}

%np. boxplot, histogramy

% \begin{figure}[H]
%     \centering
%     \includegraphics[width=0.99\linewidth]{Podstawowa_wizualizacja.png}
%     \caption{Wizualizacje}
% \end{figure}



\begin{figure}[H]
    \centering

    \begin{subfigure}[t]{0.48\linewidth}
        \centering
        \includegraphics[width=\linewidth]{podstawowa_wizualizacja_centreDistance.png}
        \caption{Wizualizacja odległości mieszkań od centrum miasta}
    \end{subfigure}
    \hfill
    \begin{subfigure}[t]{0.48\linewidth}
        \centering
        \includegraphics[width=\linewidth]{podstawowa_wizualizacja_hasBalcony.png}
        \caption{Wizualizacja ilości mieszkań z balkonem i bez balkonu}
    \end{subfigure}

    \vspace{0.5cm}

    \begin{subfigure}[t]{0.48\linewidth}
        \centering
        \includegraphics[width=\linewidth]{podstawowa_wizualizacja_poiCount.png}
        \caption{Wizualizacja ilości punktów zainteresowania w promieniu 500 metrów od mieszkań}
    \end{subfigure}
    \hfill
    \begin{subfigure}[t]{0.48\linewidth}
        \centering
        \includegraphics[width=\linewidth]{podstawowa_wizualizacja_price.png}
        \caption{Wizualizacja cen mieszkań}
    \end{subfigure}

    \vspace{0.5cm}

    \begin{subfigure}[t]{0.48\linewidth}
        \centering
        \includegraphics[width=\linewidth]{podstawowa_wizualizacja_rooms.png}
        \caption{Wizualizacja ilości pokoi w mieszkaniach}
    \end{subfigure}
    \hfill
    \begin{subfigure}[t]{0.48\linewidth}
        \centering
        \includegraphics[width=\linewidth]{podstawowa_wizualizacja_squareMeters.png}
        \caption{Wizualizacja powierzchni mieszkań}
    \end{subfigure}

    \caption{Podstawowa wizualizacja zmiennych wykorzystanych w analizie}
\end{figure}

Statystyki opisowe oraz podstawowa wizualizacja zmiennych wykorzystanych w analizie umożliwiają ocenę struktury zbioru danych, rozkładu poszczególnych cech, stopnia ich zróżnicowania oraz ewentualnych nieregularności występujących w danych.

Rozkład zmiennej \texttt{price} oraz jej podstawowe charakterystyki statystyczne wskazują na obecność obserwacji należących do różnych segmentów rynku nieruchomości. Jest to związane z tym, że w zbiorze danych występują zarówno oferty mieszkań przeznaczonych na sprzedaż, jak i oferty mieszkań oferowanych na wynajem. Ze względu na brak bezpośredniej porównywalności cen najmu i cen sprzedaży, zakres dalszej analizy ograniczmy do ofert sprzedaży mieszkań.

Dokonujemy ponownych statystyk opisowych oraz wizualizacji, uwzględniając tylko oferty sprzedaży mieszkań. Dzięki temu uzyskamy bardziej spójny i porównywalny zbiór danych, który pozwoli na rzetelną analizę skupień oraz identyfikację grup podobnych ofert w kontekście cen mieszkań przeznaczonych na sprzedaż.

\begin{table}[H]

\caption{Statystyki opisowe zmiennych wykorzystanych w analizie po usunięciu ofert wynajmu mieszkań}
\centering
\begin{tabular}[t]{lrrrrrr}
\toprule
  & Srednia & Mediana & Minimum & Maksimum & OdchylenieStandardowe & Skosnosc\\
\midrule
price & 897 888.77 & 795 000.00 & 255 000.00 & 3 000 000.00 & 394 185.67 & 1.89\\
squareMeters & 57.17 & 52.00 & 25.00 & 150.00 & 22.12 & 1.43\\
rooms & 2.63 & 3.00 & 1.00 & 6.00 & 0.95 & 0.91\\
poiCount & 23.94 & 14.00 & 0.00 & 212.00 & 33.80 & 3.07\\
centreDistance & 4.01 & 3.95 & 0.15 & 9.99 & 2.21 & 0.28\\
hasBalcony & 0.62 & 1.00 & 0.00 & 1.00 & 0.49 & -0.49\\
\bottomrule
\end{tabular}
\end{table}

\begin{figure}[H]
    \centering

    \begin{subfigure}[t]{0.48\linewidth}
        \centering
        \includegraphics[width=\linewidth]{podstawowa_wizualizacja_centreDistance_sama_sprzedaz.png}
        \caption{Wizualizacja odległości mieszkań od centrum miasta}
    \end{subfigure}
    \hfill
    \begin{subfigure}[t]{0.48\linewidth}
        \centering
        \includegraphics[width=\linewidth]{podstawowa_wizualizacja_hasBalcony_sama_sprzedaz.png}
        \caption{Wizualizacja ilości mieszkań z balkonem i bez balkonu}
    \end{subfigure}

    \vspace{0.5cm}

    \begin{subfigure}[t]{0.48\linewidth}
        \centering
        \includegraphics[width=\linewidth]{podstawowa_wizualizacja_poiCount_sama_sprzedaz.png}
        \caption{Wizualizacja ilości punktów zainteresowania w promieniu 500 metrów od mieszkań}
    \end{subfigure}
    \hfill
    \begin{subfigure}[t]{0.48\linewidth}
        \centering
        \includegraphics[width=\linewidth]{podstawowa_wizualizacja_price_sama_sprzedaz.png}
        \caption{Wizualizacja cen mieszkań}
    \end{subfigure}

    \vspace{0.5cm}

    \begin{subfigure}[t]{0.48\linewidth}
        \centering
        \includegraphics[width=\linewidth]{podstawowa_wizualizacja_rooms_sama_sprzedaz.png}
        \caption{Wizualizacja ilości pokoi w mieszkaniach}
    \end{subfigure}
    \hfill
    \begin{subfigure}[t]{0.48\linewidth}
        \centering
        \includegraphics[width=\linewidth]{podstawowa_wizualizacja_squareMeters_sama_sprzedaz.png}
        \caption{Wizualizacja powierzchni mieszkań}
    \end{subfigure}

    \caption{Podstawowa wizualizacja zmiennych wykorzystanych w analizie po usunięciu ofert wynajmu mieszkań.}
\end{figure}



% \subsection{Standaryzacja danych}


% W toku wstępnej analizy zbioru danych zauważono, że zmienna \texttt{price} nie była zapisana w jednolity sposób dla wszystkich obserwacji. Część ofert zawierała cenę całkowitą mieszkania, natomiast w przypadku części ogłoszeń cena była wyrażona jako cena za metr kwadratowy. Taka niejednorodność zapisu prowadzi do braku porównywalności obserwacji i może istotnie zaburzać wyniki dalszej analizy statystycznej oraz analizy skupień.

% Przesłanką wskazującą na występowanie dwóch różnych sposobów zapisu ceny był wyraźny uskok w rozkładzie wartości zmiennej \texttt{price}, widoczny pomiędzy poziomem około 10\,000 a 255\,000. Wartości z dolnego zakresu odpowiadają realistycznym cenom za metr kwadratowy mieszkania, natomiast wartości wyraźnie wyższe interpretowano jako ceny całkowite nieruchomości. Taki rozkład sugerował, że w zbiorze jednocześnie występują dwie różne jednostki opisu tej samej cechy.

% \begin{figure}[H]
%     \centering
%     \includegraphics[width=0.99\linewidth]{histogram_price_z_uskokiem.png}
%     \caption{Dokładna wizualizacja rozkładu zmiennej \texttt{price} z zaznaczeniem uskoku wskazującego na niejednorodność zapisu cen.}
% \end{figure}


% \begin{table}[H]
% \centering
% \begin{tabular}[t]{lrrrrrr}
% \toprule
%   & price & squareMeters & rooms & poiCount & centreDistance & hasBalcony\\
% \midrule
% \multicolumn{7}{c}{\dots}\\
% 213569 & 10 000 & 108.00 & 4 & 19 & 1.64 & 0\\
% 213620 & 10 000 & 108.00 & 4 & 19 & 1.78 & 1\\
% 214046 & 10 000 & 107.74 & 5 & 21 & 1.24 & 0\\
% 214181 & 10 000 & 110.00 & 5 & 19 & 1.28 & 1\\
% 214699 & 10 000 & 100.00 & 4 & 66 & 1.40 & 1\\
% 214804 & 10 000 & 135.00 & 4 & 72 & 1.04 & 1\\
% %\addlinespace
% 3336 & 255 000 & 37.20 & 2 & 6 & 4.13 & 0\\
% 37236 & 255 129 & 32.94 & 2 & 24 & 7.27 & 0\\
% 3782 & 266 000 & 35.50 & 2 & 9 & 4.06 & 0\\
% 20941 & 280 000 & 29.00 & 2 & 18 & 6.71 & 0\\
% 55548 & 282 000 & 41.00 & 2 & 3 & 5.44 & 1\\
% \multicolumn{7}{c}{\dots}\\
% \bottomrule
% \end{tabular}
% \caption{Obserwacje bezpośrednio przed i po uskoku w rozkładzie zmiennej \texttt{price}, ilustrujące różnicę w zapisie cen mieszkań.}
% \end{table}

% W celu ujednolicenia danych zdecydowano się na sprowadzenie wszystkich obserwacji do wspólnej postaci, to znaczy do \textbf{ceny za metr kwadratowy}. Dla obserwacji, w których cena była zapisana jako cena całkowita mieszkania, dokonano przeliczenia przez podzielenie wartości zmiennej \texttt{price} przez zmienną \texttt{squareMeters}. Obserwacje, których wartości wskazywały na to, że cena była już wyrażona w przeliczeniu na metr kwadratowy, pozostawiono bez zmian.

% Przyjęte podejście pozwoliło na usunięcie niejednorodności jednostek pomiaru i zapewniło porównywalność danych w dalszych etapach analizy. Jest to szczególnie istotne w analizie skupień, w której odległości pomiędzy obiektami są wyznaczane bezpośrednio na podstawie wartości zmiennych. Pozostawienie zmiennej \texttt{price} w niejednolitym formacie mogłoby prowadzić do sztucznego różnicowania mieszkań wyłącznie z powodu odmiennego sposobu zapisu ceny, a nie rzeczywistych różnic pomiędzy nieruchomościami.


% \begin{figure}[H]
%     \centering
%     \includegraphics[width=0.99\linewidth]{histogram_price_z_podzialem.png}
%     \caption{}
% \end{figure}


\subsection{Braki Danych}

% czy występują i jak je obsłużono

Ocena kompletności danych stanowi istotny element wstępnej analizy zbioru, ponieważ obecność brakujących wartości może wpływać na wiarygodność wyników oraz ograniczać możliwość zastosowania wybranych metod statystycznych. Z tego względu sprawdzeniu poddano wszystkie zmienne uwzględnione w dalszej analizie, tj. \texttt{price}, \texttt{squareMeters}, \texttt{rooms}, \texttt{poiCount}, \texttt{centreDistance} oraz \texttt{hasBalcony}.

W analizowanych danych nie występują braki wartości dla żadnej z uwzględnionych zmiennych. Każda obserwacja zawiera zatem pełny zestaw informacji wymaganych w dalszych etapach badania. Dzięki temu dalsza analiza nie wymaga ani uzupełniania danych, ani ograniczania zbioru do obserwacji kompletnych.

\subsection{Obserwacje odstające}

Analiza obserwacji odstających (ang. \textit{outliers}) jest kluczowym etapem przygotowania zbioru danych, ponieważ algorytmy analizy skupień są szczególnie wrażliwe na wartości ekstremalne. Mogą one prowadzić do błędnej interpretacji wyników, przesuwając środki ciężkości tworzonych grup lub tworząc sztuczne, mało liczebne skupiska.

Wartości odstające zidentyfikowano przy użyciu metody rozstępu międzykwartylowego (IQR), definiując obserwacje jako odstające, jeśli znajdowały się poza przedziałem \([Q1 - 1.5 \cdot IQR,\; Q3 + 1.5 \cdot IQR]\). Proces czyszczenia danych skupił się na następujących aspektach:

\begin{itemize}
    \item \textbf{Cena i powierzchnia:} Usunięto oferty, w których cena lub metraż przyjmowały wartości nierealne z punktu widzenia rynku nieruchomości (np. mieszkania o powierzchni powyżej 200 $m^2$ lub ceny jednostkowe drastycznie odbiegające od średniej rynkowej dla Krakowa). Wartości te mogły wynikać z błędów przy wprowadzaniu ogłoszeń.
    \item \textbf{Liczba pokoi:} Zweryfikowano rekordy, w których liczba pokoi była nieproporcjonalna do metrażu. Skupiono się na standardowych lokalach mieszkalnych, eliminując skrajne przypadki, które mogłyby zaburzyć strukturę budowanych skupień.
    \item \textbf{Dystans od centrum:} Sprawdzono poprawność zmiennej \texttt{centreDistance}, aby wykluczyć błędy w lokalizacji, które mogłyby sugerować położenie nieruchomości poza granicami badanego obszaru.
\end{itemize}
W wyniku zastosowanej procedury czyszczenia danych liczba obserwacji została zredukowana z 29\,026 do 24\,353, co stanowi około 83.9\% pierwotnego zbioru danych. Oznacza to, że usunięto około 16.1\% rekordów jako zawierających wartości odstające w co najmniej jednej z analizowanych zmiennych.

Ostatecznie, po usunięciu najbardziej skrajnych obserwacji, zbiór danych stał się bardziej spójny. Pozwoliło to na uzyskanie stabilniejszych wyników w dalszych etapach modelowania, zapewniając, że wyodrębnione skupienia będą reprezentować rzeczywiste segmenty rynku, a nie błędy w danych lub przypadki jednostkowe.



\begin{table}[H]

\caption{Statystyki opisowe po usunięciu obserwacji odstających}
\centering
\begin{tabular}[t]{lrrrrrr}
\toprule
  & Srednia & Mediana & Minimum & Maksimum & OdchylenieStandardowe & Skosnosc\\
\midrule
price & 802 044.61 & 750 000.00 & 255 000.00 & 1 649 000.00 & 250 133.54 & 0.90\\
squareMeters & 52.56 & 50.75 & 25.00 & 102.72 & 15.35 & 0.63\\
rooms & 2.48 & 2.00 & 1.00 & 4.00 & 0.76 & 0.18\\
poiCount & 14.16 & 12.00 & 0.00 & 49.00 & 10.39 & 0.88\\
centreDistance & 4.42 & 4.36 & 0.42 & 9.99 & 2.06 & 0.25\\
\addlinespace
hasBalcony & 0.64 & 1.00 & 0.00 & 1.00 & 0.48 & -0.60\\
\bottomrule
\end{tabular}
\end{table}

\begin{figure}[H]
\centering

\begin{minipage}{0.48\textwidth}
    \centering
    \includegraphics[width=\textwidth]{boxplot_price.png}
    \caption{Cena mieszkań przed czyszczeniem}
\end{minipage}
\hfill
\begin{minipage}{0.48\textwidth}
    \centering
    \includegraphics[width=\textwidth]{boxplot_price_clean.png}
    \caption{Cena mieszkań po czyszczeniu}
\end{minipage}

\vspace{0.5cm}

\begin{minipage}{0.48\textwidth}
    \centering
    \includegraphics[width=\textwidth]{boxplot_squareMeters.png}
    \caption{Metraż przed czyszczeniem}
\end{minipage}
\hfill
\begin{minipage}{0.48\textwidth}
    \centering
    \includegraphics[width=\textwidth]{boxplot_squareMeters_clean.png}
    \caption{Metraż po czyszczeniu}
\end{minipage}

\vspace{0.5cm}

\begin{minipage}{0.48\textwidth}
    \centering
    \includegraphics[width=\textwidth]{boxplot_rooms.png}
    \caption{Pokoje przed czyszczeniem}
\end{minipage}
\hfill
\begin{minipage}{0.48\textwidth}
    \centering
    \includegraphics[width=\textwidth]{boxplot_rooms_clean.png}
    \caption{Pokoje po czyszczeniu}
\end{minipage}

\end{figure}


\begin{figure}[H]
\centering

\begin{minipage}{0.48\textwidth}
    \centering
    \includegraphics[width=\textwidth]{boxplot_poiCount.png}
    \caption{Liczba punktów zainteresowania przed czyszczeniem}
\end{minipage}
\hfill
\begin{minipage}{0.48\textwidth}
    \centering
    \includegraphics[width=\textwidth]{boxplot_poiCoint_clean.png}
    \caption{Liczba punktów zainteresowania po czyszczeniu}
\end{minipage}

\vspace{0.5cm}

\begin{minipage}{0.48\textwidth}
    \centering
    \includegraphics[width=\textwidth]{boxplot_centreDistance.png}
    \caption{Odległość od centrum przed czyszczeniem}
\end{minipage}
\hfill
\begin{minipage}{0.48\textwidth}
    \centering
    \includegraphics[width=\textwidth]{boxplot_centreDistance_clean.png}
    \caption{Odległość od centrum po czyszczeniu}
\end{minipage}

\end{figure}
\chapter{Opis metod}

W niniejszym rozdziale przedstawiono dwie metody analizy skupień zastosowane w ramach projektu, tj. algorytm K-średnich (K-means) oraz hierarchiczną analizę skupień metodą aglomeracyjną. Metody te należą do najczęściej stosowanych technik uczenia nienadzorowanego i umożliwiają grupowanie obiektów na podstawie ich podobieństwa mierzonego w przestrzeni wielowymiarowej.

\section{Algorytm K-means}

\subsection{Opis metody}

Algorytm K-means jest jedną z najczęściej stosowanych metod analizy skupień \cite{macqueen1967,jain1999,hastie2009}, której celem jest podział zbioru danych na \(k\) rozłącznych klastrów w taki sposób, aby minimalizować zróżnicowanie wewnątrz poszczególnych grup.

Dla zbioru obserwacji \(X = \{x_1, x_2, \dots, x_n\}\) algorytm minimalizuje funkcję:

\[
J = \sum_{i=1}^{k} \sum_{x_j \in C_i} \|x_j - \mu_i\|^2
\]

gdzie:
\begin{itemize}
    \item \(k\) -- liczba klastrów,
    \item \(C_i\) -- i-ty klaster,
    \item \(x_j\) -- obserwacja,
    \item \(\mu_i\) -- centroid klastra,
    \item \(\|\cdot\|\) -- odległość euklidesowa.
\end{itemize}

Algorytm działa iteracyjnie, naprzemiennie przypisując obserwacje do najbliższych centroidów oraz aktualizując ich położenie jako średnią arytmetyczną przypisanych punktów. Proces ten powtarzany jest do momentu osiągnięcia zbieżności.

\subsection{Interpretacja klastrów}

Wyniki algorytmu K-means polegają na przypisaniu każdej obserwacji do jednego z \(k\) klastrów. Każdy klaster reprezentowany jest przez centroid, który można interpretować jako „typową” obserwację w danej grupie.

W kontekście analizy rynku nieruchomości klastry można interpretować jako segmenty mieszkań o podobnych cechach, takich jak cena, powierzchnia, liczba pokoi czy lokalizacja.

Interpretacja klastrów pozwala na identyfikację różnych segmentów rynku, np. mieszkań:
\begin{itemize}
    \item o wysokim standardzie i dużej powierzchni,
    \item średniego segmentu cenowego,
    \item małych mieszkań o niższej cenie i gorszej lokalizacji.
\end{itemize}

Dzięki temu możliwe jest nie tylko grupowanie danych, ale również ich praktyczna interpretacja ekonomiczna, co stanowi główną wartość analizy skupień w niniejszym projekcie.

\subsection{Dobór liczby klastrów}

Jednym z kluczowych etapów algorytmu K-means jest wybór odpowiedniej liczby klastrów \(k\), ponieważ metoda ta nie wyznacza tej wartości automatycznie.

W niniejszej pracy zastosowano metodę tzw. \textit{elbow method} (metody łokcia), która polega na analizie wartości funkcji celu (sumy kwadratów odległości obserwacji od centroidów) dla różnych wartości \(k\).

Dla każdej liczby klastrów obliczono miarę:

\[
J(k) = \sum_{i=1}^{k} \sum_{x_j \in C_i} \|x_j - \mu_i\|^2
\]

Następnie analizowano, jak zmienia się wartość \(J(k)\) wraz ze wzrostem liczby klastrów.

Wraz ze wzrostem \(k\) wartość błędu maleje, jednak po pewnym momencie tempo tej poprawy znacząco spada. Punkt ten, określany jako \textit{łokieć}, uznaje się za optymalną liczbę klastrów, ponieważ dalsze zwiększanie liczby grup nie przynosi istotnej poprawy jakości podziału danych.

Zastosowanie tej metody pozwala na kompromis pomiędzy jakością grupowania a jego złożonością interpretacyjną.

\section{Analiza hierarchiczna klastrów}

\subsection{Opis metody}

Analiza hierarchiczna klastrów (ang. \textit{hierarchical clustering}) jest metodą grupowania danych, która nie wymaga wcześniejszego określenia liczby klastrów. Zamiast tego buduje ona hierarchię skupień w postaci drzewa zależności, tzw. dendrogramu.

Metoda ta rozpoczyna się od traktowania każdej obserwacji jako osobnego klastra, a następnie w kolejnych krokach łączy najbardziej podobne grupy, aż do uzyskania jednego wspólnego klastra obejmującego wszystkie obserwacje.

Do określenia podobieństwa pomiędzy obserwacjami wykorzystywana jest miara odległości, najczęściej odległość euklidesowa:

\[
d(x_i, x_j) = \sqrt{\sum_{p=1}^{m} (x_{ip} - x_{jp})^2}
\]

gdzie:
\begin{itemize}
    \item \(x_i, x_j\) -- obserwacje,
    \item \(m\) -- liczba zmiennych,
    \item \(d(x_i, x_j)\) -- odległość pomiędzy obserwacjami.
\end{itemize}

W procesie łączenia klastrów istotną rolę odgrywa metoda wiązania (ang. \textit{linkage}), która określa sposób obliczania odległości pomiędzy grupami. W niniejszej pracy zastosowano metodę wiązania Warda \cite{ward1963}, która minimalizuje wzrost wariancji wewnątrz klastrów przy każdym połączeniu.

\subsection{Interpretacja metody}

Wynikiem analizy hierarchicznej jest dendrogram, czyli drzewo przedstawiające kolejność łączenia obserwacji i klastrów.

Przecięcie dendrogramu na wybranym poziomie pozwala uzyskać określoną liczbę klastrów, które można interpretować jako segmenty danych o podobnych cechach.

W kontekście analizy rynku nieruchomości metoda ta pozwala na identyfikację naturalnej struktury danych oraz ocenę, jak poszczególne grupy mieszkań łączą się ze sobą w zależności od ich podobieństwa.

Dodatkową zaletą analizy hierarchicznej jest możliwość wizualnej oceny relacji pomiędzy obserwacjami, co ułatwia interpretację struktury danych oraz wybór liczby klastrów.

\section{Podsumowanie metod}

Zastosowanie dwóch różnych metod analizy skupień pozwala na porównanie stabilności oraz jakości otrzymanych wyników. K-means zapewnia efektywność obliczeniową i prostotę implementacji, natomiast metoda hierarchiczna umożliwia bardziej szczegółową analizę struktury danych i wybór liczby klastrów na podstawie dendrogramu.

\section{Przegląd literatury dotyczącej metod}

Metoda K-means została po raz pierwszy formalnie opisana przez Stuarta Lloyda w 1957 roku w kontekście kompresji sygnałów, a następnie szeroko rozpowszechniona w literaturze statystycznej i uczeniu maszynowym.

Jedną z najczęściej cytowanych prac dotyczących algorytmu K-means jest publikacja MacQueena, który w 1967 roku wprowadził pojęcie klastrów w kontekście analizy danych:

\begin{itemize}
    \item MacQueen, J. (1967). \textit{Some Methods for classification and Analysis of Multivariate Observations}. Proceedings of the Fifth Berkeley Symposium on Mathematical Statistics and Probability.
\end{itemize}

Współczesne ujęcie metod analizy skupień, w tym zarówno metody K-means, jak i metod hierarchicznych, zostało szczegółowo opisane w literaturze podręcznikowej dotyczącej eksploracyjnej analizy danych, m.in.:

\begin{itemize}
    \item Jain, A. K., Murty, M. N., \& Flynn, P. J. (1999). \textit{Data Clustering: A Review}. ACM Computing Surveys.
    \item Hastie, T., Tibshirani, R., \& Friedman, J. (2009). \textit{The Elements of Statistical Learning}. Springer.
\end{itemize}

Prace te stanowią podstawę teoretyczną dla zastosowanych w niniejszym projekcie metod analizy skupień oraz ich interpretacji w kontekście danych wielowymiarowych \cite{macqueen1967,ward1963,jain1999,hastie2009}.







\chapter{Wyniki}

W niniejszym rozdziale zaprezentowano wyniki przeprowadzonej analizy skupień dla rynku nieruchomości w Krakowie. Proces badawczy obejmował wyznaczenie optymalnej liczby klastrów, podział zbioru danych algorytmem K-means oraz weryfikację struktury danych metodą hierarchiczną.

\section{Wybór optymalnej liczby skupień}

Pierwszym etapem analizy było wyznaczenie liczby skupień ($k$) przy użyciu metody łokcia (ang. \textit{elbow method}). Analizie poddano sumę kwadratów odległości wewnątrz skupień dla wartości $k$ z przedziału od 1 do 10. Wyniki przedstawiono na rysunku \ref{fig:elbow}.

\begin{figure}[H]
    \centering
    \includegraphics[width=0.75\linewidth]{metoda_lokcia.png}
    \caption{Wykres metody łokcia dla danych krakowskich ofert mieszkań}
    \label{fig:elbow}
\end{figure}

Na powyższym wykresie punkt zagięcia krzywej (tzw. łokieć) jest wyraźnie widoczny dla $k=3$. Oznacza to, że dalsze zwiększanie liczby grup nie powoduje już tak znaczącego spadku zmienności wewnątrz skupień. W związku z tym, do dalszych badań przyjęto podział na trzy segmenty.

\section{Charakterystyka klastrów (K-means)}

Po zastosowaniu algorytmu K-means dla $k=3$ na ustandaryzowanych danych, otrzymano podział zbioru na trzy grupy o zróżnicowanej charakterystyce. W tabeli \ref{tab:klastry} zestawiono średnie wartości zmiennych dla każdego z segmentów (w oryginalnych jednostkach).

\begin{table}[H]
\caption{Średnie wartości zmiennych w wyodrębnionych skupieniach}
\label{tab:klastry}
\centering
\begin{tabular}{lrrr}
\toprule
Zmienna & Klaster 1 & Klaster 2 & Klaster 3 \\
\midrule
Cena [zł] & 1 657 913,56 & 783 671,32 & 1 120 747,56 \\
Powierzchnia [$m^2$] & 94,84 & 45,17 & 64,81 \\
Liczba pokoi & 3,74 & 2,05 & 2,86 \\
Liczba POI w okolicy & 16,92 & 13,87 & 13,53 \\
Odległość od centrum [km] & 5,86 & 5,34 & 10,75 \\
Balkon (1=tak) & 0,76 & 0,38 & 0,79 \\
\bottomrule
\end{tabular}
\end{table}

\section{Interpretacja segmentów rynku}

Na podstawie analizy statystycznej centrów skupień można zdefiniować następujące typy ofert występujących na krakowskim rynku:

\begin{itemize}
    \item \textbf{Skupienie 1 (Segment Premium/Duże Metraże):} Grupa skupiająca najdroższe oferty o największej powierzchni oraz największej liczbie pokoi. Są to często apartamenty w prestiżowych lokalizacjach lub duże mieszkania rodzinne.
    
    \item \textbf{Skupienie 2 (Segment Kompaktowy/Ekonomiczny):} Obejmuje mieszkania o najniższej średniej cenie, najmniejszej powierzchni i najmniejszej liczbie pokoi. Są to zazwyczaj mniejsze lokale, bardziej dostępne cenowo, skierowane m.in. do singli, par lub inwestorów.
    
    \item \textbf{Skupienie 3 (Segment Peryferyjny/Pośredni):} Charakteryzuje się pośrednimi wartościami ceny, metrażu i liczby pokoi, ale największą odległością od centrum miasta. Są to zazwyczaj mieszkania położone na obrzeżach miasta, stanowiące kompromis między parametrami użytkowymi a lokalizacją.
\end{itemize}

\section{Weryfikacja metodą hierarchiczną}

W celu potwierdzenia istnienia naturalnych skupisk w danych, przeprowadzono klastrowanie hierarchiczne metodą Warda na reprezentatywnej próbie obserwacji. Dendrogram (Rys. \ref{fig:dendrogram}) potwierdził zasadność podziału na trzy główne gałęzie, co jest spójne z wynikami metody K-means.

\begin{figure}[H]
    \centering
    \includegraphics[width=0.85\linewidth]{dendrogram.png}
    \caption{Dendrogram analizy hierarchicznej (próba $n=100$)}
    \label{fig:dendrogram}
\end{figure}

\section{Porównanie uzyskanych wyników i ocena jakości skupień}

Dokonano porównania wyników uzyskanych dwiema niezależnymi metodami: algorytmem K-means oraz metodą hierarchiczną Warda. 

\begin{itemize}
    \item \textbf{Zbieżność podziału:} Obie zastosowane metody wskazały na istnienie trzech dominujących skupień w badanym zbiorze danych. Metoda łokcia (Rys. \ref{fig:elbow}) oraz struktura dendrogramu (Rys. \ref{fig:dendrogram}) wykazują wysoką spójność, co świadczy o stabilności wyodrębnionych segmentów rynku.
    \item \textbf{Ocena jakości:} Skupienia charakteryzują się wysoką jednorodnością wewnątrzgałęziową. Klaster 1 (Premium), Klaster 2 (Kompaktowy/Ekonomiczny) oraz Klaster 3 (Peryferyjny/Pośredni) różnią się wyraźnie pod względem ceny, metrażu i lokalizacji, co wskazuje na sensowność otrzymanego podziału rynku.
    \item \textbf{Weryfikacja założeń:} Wykorzystanie ustandaryzowanych zmiennych pozwoliło na uniknięcie dominacji zmiennej „Cena” nad pozostałymi deskryptorami, co przekłada się na merytoryczną poprawność klastrowania.
\end{itemize}

Podsumowując, zastosowanie dwóch różnych podejść algorytmicznych pozwoliło na uzyskanie wiarygodnej segmentacji krakowskiego rynku nieruchomości, co realizuje główny cel badania.

\chapter{Podsumowanie}

%(ocena realizacji celu, odniesienie do pozycji z przeglądu literatury)


Celem niniejszej pracy była identyfikacja jednorodnych grup ofert sprzedaży mieszkań w Krakowie przy wykorzystaniu metod analizy skupień. Założono, że na podstawie wybranych cech nieruchomości, takich jak cena, powierzchnia, liczba pokoi, odległość od centrum, liczba punktów zainteresowania w otoczeniu oraz informacja o posiadaniu balkonu, możliwe będzie wyodrębnienie segmentów rynku różniących się charakterystyką i poziomem atrakcyjności. Na podstawie przeprowadzonych analiz można stwierdzić, że cel pracy został zrealizowany.

\section{Ocena realizacji celu}

\textbf{Po pierwsze, w danych można wskazać wyraźne skupienia ofert mieszkań.} Zastosowanie metody łokcia w analizie algorytmu K-means wskazało, że optymalnym rozwiązaniem jest podział zbioru na trzy klastry. Wynik ten został dodatkowo potwierdzony przez analizę hierarchiczną metodą Warda, której dendrogram również sugerował istnienie trzech głównych grup. Oznacza to, że badany zbiór nie ma charakteru przypadkowego ani jednorodnego, lecz wykazuje wewnętrzną strukturę, którą można opisać za pomocą metod grupowania.

\textbf{Po drugie, wyodrębnione klastry dają się sensownie zinterpretować ekonomicznie.} Otrzymane skupienia różniły się między sobą przede wszystkim poziomem cen, powierzchnią mieszkań, liczbą pokoi oraz położeniem względem centrum miasta. Pozwala to traktować je jako odmienne segmenty rynku mieszkaniowego, obejmujące oferty o różnym standardzie, funkcji i dostępności cenowej. Tym samym analiza skupień nie doprowadziła jedynie do technicznego podziału danych, lecz umożliwiła także sformułowanie wniosków dotyczących struktury rynku nieruchomości w Krakowie.

\textbf{Po trzecie, wyniki uzyskane różnymi metodami analizy skupień są ze sobą zgodne.} Zarówno algorytm K-means, jak i hierarchiczna analiza skupień prowadzą do wniosku, że najbardziej uzasadniony jest podział na trzy główne segmenty rynku. Zbieżność wyników uzyskanych niezależnymi metodami zwiększa wiarygodność przeprowadzonej analizy i wskazuje, że otrzymany podział nie jest efektem przypadku ani specyfiki pojedynczego algorytmu.

\section{Ocena uzyskanych wyników w kontekście literatury przedmiotu}

Uzyskane wyniki są zgodne z wnioskami przedstawionymi w literaturze przedmiotu. Bourassa, Hamelink, Hoesli i MacGregor \cite{bourassa1999} wskazywali, że subrynki mieszkaniowe mogą być wyodrębniane empirycznie przy użyciu metod statystycznych, w tym analizy skupień, a taki podział może lepiej odzwierciedlać zróżnicowanie rynku niż klasyfikacje arbitralne. Wyniki niniejszej pracy potwierdzają tę obserwację, ponieważ zastosowane metody pozwoliły wyodrębnić sensowne segmenty rynku mieszkań na podstawie rzeczywistych cech ofert.

Przeprowadzona analiza pozostaje również spójna z badaniem Tomala \cite{tomal2021}, który wykazał, że rynek mieszkaniowy w Polsce cechuje się wyraźną heterogenicznością, a zastosowanie metod klasteryzacji umożliwia identyfikację jednorodnych grup rynków różniących się poziomem cen i cechami zasobu mieszkaniowego. W niniejszej pracy podobny mechanizm zaobserwowano na niższym poziomie agregacji, tj. w obrębie jednego miasta. Oznacza to, że zróżnicowanie rynku mieszkaniowego jest widoczne nie tylko pomiędzy rynkami lokalnymi w Polsce, ale także wewnątrz pojedynczego rynku miejskiego.

Wyniki badania można także odnieść do pracy Tomala i Helbicha \cite{tomalhelbich2022}, którzy pokazali, że rynek mieszkaniowy w Krakowie ma wewnętrzną strukturę i może być analizowany jako zbiór subrynków o odmiennej charakterystyce. Choć ich badanie dotyczyło rynku najmu, a nie sprzedaży, uzyskane przez nich wnioski o niejednorodności rynku krakowskiego znajdują potwierdzenie również w niniejszej pracy. Otrzymane skupienia wskazują bowiem, że oferty mieszkań w Krakowie różnią się nie tylko ceną, ale także relacją między wielkością lokalu, jego położeniem oraz cechami otoczenia.

Należy jednocześnie zauważyć, że przeprowadzona analiza ma pewne ograniczenia. Badanie oparto na jednym mieście oraz na ograniczonym zestawie zmiennych, co oznacza, że uzyskane klastry nie wyczerpują całej złożoności rynku nieruchomości. W dalszych pracach możliwe byłoby rozszerzenie analizy o dodatkowe cechy, takie jak standard wykończenia, rok budowy, piętro czy forma własności, a także porównanie wyników dla innych miast. Mimo tych ograniczeń uzyskane rezultaty pozwalają uznać, że analiza skupień stanowi użyteczne narzędzie segmentacji rynku mieszkaniowego.

Podsumowując, cel pracy został osiągnięty: w danych dotyczących ofert sprzedaży mieszkań w Krakowie zidentyfikowano wyraźne i interpretowalne skupienia, a otrzymane wyniki okazały się spójne zarówno pomiędzy zastosowanymi metodami, jak i z ustaleniami przedstawionymi w literaturze przedmiotu.























\begin{thebibliography}{9}

\bibitem{bourassa1999}
Bourassa, S. C., Hamelink, F., Hoesli, M., MacGregor, B. D. (1999).
\textit{Defining Housing Submarkets}.
Journal of Housing Economics, 8(2), 160--183.
\url{https://doi.org/10.1006/jhec.1999.0246}

\bibitem{tomal2021}
Tomal, M. (2021).
\textit{Housing market heterogeneity and cluster formation: evidence from Poland}.
International Journal of Housing Markets and Analysis, 14(5), 1166--1185.
\url{https://doi.org/10.1108/IJHMA-09-2020-0114}

\bibitem{tomalhelbich2022}
Tomal, M., Helbich, M. (2022).
\textit{The private rental housing market before and during the COVID-19 pandemic: A submarket analysis in Cracow, Poland}.
Environment and Planning B: Urban Analytics and City Science, 49(6), 1646--1662.
\url{https://doi.org/10.1177/23998083211062907}

\bibitem{macqueen1967}
MacQueen, J. B. (1967).
\textit{Some Methods for Classification and Analysis of Multivariate Observations}.
In: Le Cam, L. M., Neyman, J. (red.),
\textit{Proceedings of the Fifth Berkeley Symposium on Mathematical Statistics and Probability},
Vol. 1, University of California Press, 281--297.

\bibitem{ward1963}
Ward, J. H. Jr. (1963).
\textit{Hierarchical Grouping to Optimize an Objective Function}.
Journal of the American Statistical Association, 58(301), 236--244.
\url{https://doi.org/10.1080/01621459.1963.10500845}

\bibitem{jain1999}
Jain, A. K., Murty, M. N., Flynn, P. J. (1999).
\textit{Data Clustering: A Review}.
ACM Computing Surveys, 31(3), 264--323.
\url{https://doi.org/10.1145/331499.331504}

\bibitem{hastie2009}
Hastie, T., Tibshirani, R., Friedman, J. (2009).
\textit{The Elements of Statistical Learning: Data Mining, Inference, and Prediction}.
2nd ed., Springer.
\url{https://doi.org/10.1007/978-0-387-84858-7}

\end{thebibliography}

\end{document}
